\documentclass[11pt]{article}
\usepackage[margin=1in]{geometry}
\usepackage[T1]{fontenc}
\usepackage{lmodern}
\usepackage{amsmath,amssymb,bm,enumitem}
\usepackage[nopatch=footnote]{microtype}
\usepackage{hyperref}
\hypersetup{colorlinks=true,linkcolor=blue,citecolor=blue,urlcolor=blue}
\setlength{\parindent}{0pt}
\setlength{\parskip}{0.45em}
\setlist{nosep,leftmargin=2em}
\newcommand{\by}{\bm{y}}
\newcommand{\bx}{\bm{x}}
\newcommand{\bo}{\bm{o}}
\newcommand{\bw}{\bm{w}}
\newcommand{\bu}{\bm{u}}
\newcommand{\bv}{\bm{v}}
\newcommand{\br}{\bm{r}}
\newcommand{\bb}{\bm{b}}
\newcommand{\bz}{\bm{z}}
\newcommand{\balpha}{\bm{\alpha}}
\newcommand{\bbeta}{\bm{\beta}}
\newcommand{\btheta}{\bm{\theta}}
\newcommand{\bxi}{\bm{\xi}}
\newcommand{\bmu}{\bm{\mu}}
\newcommand{\bA}{\bm{A}}
\newcommand{\bB}{\bm{B}}
\newcommand{\bC}{\bm{C}}
\newcommand{\bD}{\bm{D}}
\newcommand{\bH}{\bm{H}}
\newcommand{\bI}{\bm{I}}
\newcommand{\bK}{\bm{K}}
\newcommand{\bL}{\bm{L}}
\newcommand{\bM}{\bm{M}}
\newcommand{\bP}{\bm{P}}
\newcommand{\bQ}{\bm{Q}}
\newcommand{\bR}{\bm{R}}
\newcommand{\bS}{\bm{S}}
\newcommand{\bT}{\bm{T}}
\newcommand{\bV}{\bm{V}}
\newcommand{\bW}{\bm{W}}
\newcommand{\bX}{\bm{X}}
\newcommand{\bZ}{\bm{Z}}
\newcommand{\bSigma}{\bm{\Sigma}}
\newcommand{\bLambda}{\bm{\Lambda}}
\newcommand{\bzero}{\bm{0}}
\newcommand{\trans}{^{\mathsf T}}
\DeclareMathOperator{\diag}{diag}
\DeclareMathOperator{\blockdiag}{blockdiag}
\DeclareMathOperator{\Cov}{Cov}
\DeclareMathOperator{\tr}{tr}
\title{\texttt{stLMM}: Algebra for IID and Structured Random Effects}
\author{}
\date{Development specification --- \today}

\begin{document}
\maketitle
\begin{abstract}
This note develops a common Gaussian latent-variable formulation for grouped
IID coefficients and structured effects. It distinguishes fixed observation
designs from parameter-dependent loadings, derives separate and jointly collapsed
calculations, and states the requirements for likelihood augmentation,
multivariate processes, and spatial misalignment. The jointly IID-collapsed
sampler is a proposed extension to the version 0.0.6 baseline. The equations do
not imply that every model construction discussed here is already implemented.
\end{abstract}
{\small\tableofcontents}
\clearpage

\section{Observation model and notation}
For the $n$ observed responses, consider
\begin{equation}\label{eq:observation}
 \by = \bo + \bX\bbeta + \bZ\balpha + \bA\bw + \bm{\epsilon},
 \qquad \bm{\epsilon}\sim N_n(\bzero,\bR).
\end{equation}
Here $\bo$ is a known offset, $\bX$ is the fixed-effect design, $\bZ$ maps grouped
coefficients to observations, and $\bA$ maps structured latent coefficients to
observations. Use $\bH$ for the \emph{combined} observation mapping so that the
existing meaning of $\bA$ is preserved. The more general specification is
\begin{equation}\label{eq:general}
 \by\mid\bbeta,\bv,\btheta \sim
 N_n\{\bo+\bX\bbeta+\bH(\btheta)\bv,\bR(\btheta)\},
 \qquad
 \bv\mid\btheta\sim N_q\{\bzero,\bQ(\btheta)^{-1}\}.
\end{equation}
The parameter vector $\btheta$ includes covariance parameters and, where present,
unknown loadings. Both $\bH$ and $\bQ$ may depend on these parameters. Conditional
on their current values, their entries are numerical quantities for the Gaussian
calculations. Symbols without their arguments below retain this dependence.

The derivations assume proper positive-definite $\bQ$ and $\bR$ and set
$\bW=\bR^{-1}$. The present observation models use diagonal $\bW$; the identities
also hold for positive-definite block observation precision. Intrinsic or otherwise
singular priors need a constrained formulation and are not covered by ordinary
determinants in this note. Missing responses contribute no observation-likelihood
row; the latent support may still contain their locations or group levels.

For independent IID and structured prior blocks, the centered specification is
\begin{equation}\label{eq:stack}
 \bv=\begin{pmatrix}\balpha\\\bw\end{pmatrix},\qquad
 \bH=[\bZ\ \bA],\qquad
 \bQ=\blockdiag(\bD_\alpha^{-1},\bQ_w).
\end{equation}
The block-diagonal prior expresses independence between these components given
parameters. Their posterior distributions generally are dependent.

The representation selected for this extension keeps $\bH=[\bZ\ \bA]$ fixed
and puts covariance parameters in $\bQ$. It targets full-rank vector-valued IID
and structured terms, including multivariate responses with partial observation.
The parameter-dependent mapping identities describe alternative representations;
reduced-rank factor models and unknown observation loadings are outside this
extension's scope.

\section{Correlated IID coefficients with a fixed design}
For a grouped term with $G$ groups and $k$ coefficients per group, write
\begin{equation}\label{eq:iid}
 \balpha_g\mid\bSigma\stackrel{\mathrm{iid}}{\sim}N_k(\bzero,\bSigma),
 \qquad \bD_\alpha=\bI_G\otimes\bSigma.
\end{equation}
The IID designation applies to the coefficient \emph{vectors} across groups.
Entries within each vector may be correlated. For independent terms, concatenate
their designs and use independent prior blocks. For parent-specific covariance
matrices, replace each term block by $\blockdiag_h(\bI_{G_h}\otimes\bSigma_h)$.
These covariance specifications are distinct from whether the group memberships
are crossed or nested.

For a correlated intercept and slope,
\begin{equation}\label{eq:two-coefficients}
 \balpha_g=\begin{pmatrix}\alpha_{0g}\\\alpha_{1g}\end{pmatrix},\qquad
 \bSigma=\begin{pmatrix}
 \sigma_0^2 & \rho\sigma_0\sigma_1\\
 \rho\sigma_0\sigma_1 & \sigma_1^2
 \end{pmatrix},
 \qquad \sigma_0,\sigma_1>0,\quad |\rho|<1.
\end{equation}
An observation in group $g$ receives the contribution
\begin{equation}
 \alpha_{0g}+x_i\alpha_{1g}=[1\ \ x_i]\balpha_g.
\end{equation}
Its row of $\bZ$ therefore contains $[1\ \ x_i]$ in the columns for that group,
regardless of $\rho$. For observations $i,j$ in the same group,
\begin{equation}\label{eq:iid-crosscov}
 \Cov(\alpha_{0g}+x_i\alpha_{1g},\alpha_{0g}+x_j\alpha_{1g}\mid\bSigma)
 =\sigma_0^2+\rho\sigma_0\sigma_1(x_i+x_j)+\sigma_1^2x_ix_j.
\end{equation}
The off-diagonal covariance is thus present even though the design row is fixed.
For this term, effects belonging to different groups have zero prior covariance.
The prior precision contributes repeated small blocks
\begin{equation}\label{eq:two-precision}
 \bSigma^{-1}=\frac{1}{1-\rho^2}
 \begin{pmatrix}
 \sigma_0^{-2} & -\rho/(\sigma_0\sigma_1)\\
 -\rho/(\sigma_0\sigma_1) & \sigma_1^{-2}
 \end{pmatrix}.
\end{equation}
Updating $\rho$ or the variances changes these blocks of $\bQ$, not the supplied
covariates or group memberships in $\bZ$.

\subsection{Equivalent standardized coefficients and variable loadings}
Let $\bSigma=\bC\bC\trans$ and set
\begin{equation}\label{eq:standardize}
 \balpha_g=\bC\bu_g,\qquad \bu_g\stackrel{\mathrm{iid}}{\sim}N_k(\bzero,\bI_k).
\end{equation}
The design becomes $\bZ(\bI_G\otimes\bC)$. In the intercept--slope example,
\begin{equation}\label{eq:variable-row}
 \bC=\begin{pmatrix}
 \sigma_0&0\\\rho\sigma_1&\sigma_1\sqrt{1-\rho^2}
 \end{pmatrix},\qquad
 [1\ \ x_i]\bC=
 [\sigma_0+x_i\rho\sigma_1\ \ \ x_i\sigma_1\sqrt{1-\rho^2}].
\end{equation}
These mapping entries are unknown functions of covariance parameters and change
across posterior draws. In the combined system,
\begin{equation}
 \bH_{\mathrm{std}}=[\bZ(\bI_G\otimes\bC)\ \bA],\qquad
 \bQ_{\mathrm{std}}=\blockdiag(\bI_{Gk},\bQ_w).
\end{equation}
The two representations induce exactly the same observation covariance and
marginal likelihood. They define the same posterior model when the prior on
$\bSigma$ is preserved under any parameter transformation. An arbitrary new prior
on the entries of $\bC$ need not preserve that model.

More generally, if $\bv=\bT\bu$ for a nonsingular $\bT$, then
$\bH_u=\bH\bT$ and $\bQ_u=\bT\trans\bQ\bT$. Consequently,
\begin{equation}\label{eq:invariance}
 \bH_u\bQ_u^{-1}\bH_u\trans=\bH\bQ^{-1}\bH\trans.
\end{equation}
This identity holds at each parameter value even when $\bT$ depends on parameters.
Collapsing either latent representation does not itself change the marginal
target for identically parameterized covariance parameters. Numerical cost and
conditioning can differ.

\section{Jointly collapsed Gaussian calculations}
Integrating $\bv$ out of \eqref{eq:general} gives
\begin{equation}\label{eq:marginal}
 \by\mid\bbeta,\btheta\sim N_n(\bo+\bX\bbeta,\bV),\qquad
 \bV=\bR+\bH\bQ^{-1}\bH\trans.
\end{equation}
Define
\begin{equation}\label{eq:M}
 \bM=\bQ+\bH\trans\bW\bH,\qquad
 \br=\by-\bo-\bX\bbeta,\qquad \bb=\bH\trans\bW\br.
\end{equation}
Completing the square in the joint Gaussian kernel yields
\begin{align}
 &(\br-\bH\bv)\trans\bW(\br-\bH\bv)+\bv\trans\bQ\bv\notag\\
 &\quad=(\bv-\bM^{-1}\bb)\trans\bM(\bv-\bM^{-1}\bb)
       +\br\trans\bW\br-\bb\trans\bM^{-1}\bb.
\end{align}
The determinant lemma and Woodbury identity give
\begin{align}
 \log|\bV|&=\log|\bR|+\log|\bM|-\log|\bQ|,\label{eq:det}\\
 \bV^{-1}\br&=\bW\br-\bW\bH\bM^{-1}\bH\trans\bW\br,\label{eq:solve}\\
 \br\trans\bV^{-1}\br&=\br\trans\bW\br-\bb\trans\bM^{-1}\bb.\label{eq:quad}
\end{align}
No observation covariance inverse is required: each occurrence of $\bM^{-1}$
denotes a solve. The Gaussian log likelihood is
\begin{equation}\label{eq:loglik}
 \ell=-\frac n2\log(2\pi)-\frac12\left\{
 \log|\bR|+\log|\bM|-\log|\bQ|
 +\br\trans\bW\br-\bb\trans\bM^{-1}\bb\right\}.
\end{equation}
Equations \eqref{eq:M}--\eqref{eq:loglik} allow both $\bH(\btheta)$ and
$\bQ(\btheta)$. Any parameter-dependent entries must be refreshed when evaluating
a proposal. In the centered IID case,
\begin{equation}\label{eq:M-blocks}
 \bM=\begin{pmatrix}
 \bD_\alpha^{-1}+\bZ\trans\bW\bZ & \bZ\trans\bW\bA\\
 \bA\trans\bW\bZ & \bQ_w+\bA\trans\bW\bA
 \end{pmatrix},\qquad
 \log|\bQ|=\log|\bQ_w|-\sum_{b,h}G_{bh}\log|\bSigma_{bh}|.
\end{equation}
The sum denotes the appropriate independent term/parent blocks, including scalar
blocks as $1\times1$ covariance matrices.

\subsection{Fixed effects and latent recovery}
For a normal fixed-effect prior with precision $\bP_\beta$ and mean $\bmu_\beta$,
the conditional precision and mean after integrating the latent effects are
\begin{align}
 \bK_\beta&=\bP_\beta+\bX\trans\bV^{-1}\bX,\\
 \bm{m}_\beta&=\bK_\beta^{-1}
 \{\bP_\beta\bmu_\beta+\bX\trans\bV^{-1}(\by-\bo)\}.
\end{align}
A flat prior corresponds to $\bP_\beta=0$ when the resulting conditional precision
is positive definite. The latent conditional is
\begin{equation}\label{eq:latent-recovery}
 \bv\mid\by,\bbeta,\btheta\sim N_q(\bM^{-1}\bb,\bM^{-1}).
\end{equation}
If a permutation satisfies $\bP\bM\bP\trans=\bL\bL\trans$, a draw is
\begin{equation}
 \bv=\bM^{-1}\bb+\bP\trans\bL^{-\mathsf T}\bz,
 \qquad \bz\sim N_q(\bzero,\bI_q).
\end{equation}
This describes triangular solves, not an explicitly formed inverse or a requirement
to extract the factor from CHOLMOD. The IID and structured components of a joint
posterior draw generally cannot be replaced by independent draws from their two
marginals. Output storage may retain selected components without changing the
joint draw used internally.

\section{Keeping IID effects separate}
When only $\bw$ is integrated out, define
\begin{equation}
 \bV_w=\bR+\bA\bQ_w^{-1}\bA\trans,\qquad
 \bM_w=\bQ_w+\bA\trans\bW\bA.
\end{equation}
The observation mean is then $\bo+\bX\bbeta+\bZ\balpha$. Conditional on
$\bbeta$ and covariance parameters,
\begin{align}
 \bK_\alpha&=\bD_\alpha^{-1}+\bZ\trans\bV_w^{-1}\bZ,\label{eq:alpha-conditional}\\
 \bm{m}_\alpha&=\bK_\alpha^{-1}\bZ\trans\bV_w^{-1}
 (\by-\bo-\bX\bbeta).
\end{align}
The fixed-effect conditional similarly uses $\by-\bo-\bZ\balpha$ and $\bV_w$.
The Schur complement of the structured block in \eqref{eq:M-blocks} is
\begin{align}
 &\bD_\alpha^{-1}+\bZ\trans\bW\bZ
 -\bZ\trans\bW\bA\bM_w^{-1}\bA\trans\bW\bZ\notag\\
 &\qquad=\bD_\alpha^{-1}+\bZ\trans\bV_w^{-1}\bZ
 =\bK_\alpha.\label{eq:schur}
\end{align}
Thus the two systems are algebraically connected. Keeping IID coefficients
separate avoids adding their nodes to the structured factorization, but its
conditional IID precision can be dense.

For one covariance block, let $\bS_\alpha=\sum_{g=1}^G\balpha_g\balpha_g\trans$.
Its covariance-parameter target conditional on the coefficients is
\begin{equation}\label{eq:conditional-covariance-target}
 \log p(\bSigma\mid\balpha,\ldots)
 =\log p(\bSigma)-\frac G2\log|\bSigma|
 -\frac12\tr(\bSigma^{-1}\bS_\alpha)+\mathrm{constant}.
\end{equation}
For the inverse-Wishart convention
\begin{equation}
 p(\bSigma)\propto |\bSigma|^{-(\nu+k+1)/2}
 \exp\{-\tfrac12\tr(\bm{\Psi}\bSigma^{-1})\},
\end{equation}
with $\nu>k-1$ and positive-definite $\bm{\Psi}$, the conditional is
$\operatorname{IW}_k(\nu+G,\bm{\Psi}+\bS_\alpha)$.
For a scalar variance with $\operatorname{IG}(a,b)$ density proportional to
$v^{-a-1}\exp(-b/v)$, it is
$\operatorname{IG}(a+G/2,b+\sum_g\alpha_g^2/2)$.
Other priors can use \eqref{eq:conditional-covariance-target} in a Metropolis step.
That step uses small covariance blocks and does not require changing $\bM_w$.

\section{Marginal covariance targets and prior transformations}
When the IID effects are integrated out in the covariance update, the target is
\begin{equation}\label{eq:marginal-target}
 \log p(\btheta\mid\by,\bbeta)=\ell(\btheta)+\log p(\btheta)+\mathrm{constant}.
\end{equation}
The IG/IW priors remain valid choices, but their usual conditional conjugacy is
generally lost. For an unconstrained proposal coordinate $\bxi$ with
$\btheta=t(\bxi)$, the target density with respect to $\bxi$ is
\begin{equation}
 \ell\{t(\bxi)\}+\log p\{t(\bxi)\}
 +\log\left|\frac{\partial t(\bxi)}{\partial\bxi}\right|.
\end{equation}
The prior measure must be specified. A prior on a full covariance matrix differs
from independent priors on its marginal variances and correlation matrix.

\subsection{Log-Cholesky proposals for an IW covariance}
The selected IW proposal coordinates use $\bSigma=\bL\bL\trans$ with a lower
triangular $\bL$, positive diagonal, $\xi_{ii}=\log L_{ii}$, and
$\xi_{ij}=L_{ij}$ for $i>j$. The Jacobian from the distinct entries of
$\bL$ to the distinct covariance entries is \cite{Banerjee2008}
\begin{equation}
 J_{\Sigma\leftarrow L}=2^k\prod_{i=1}^k L_{ii}^{\,k-i+1}.
\end{equation}
Exponentiating the diagonal contributes another $\prod_i L_{ii}$. Thus
\begin{equation}\label{eq:log-cholesky-jacobian}
 \log J_{\Sigma\leftarrow\xi}
 =k\log 2+\sum_{i=1}^k(k-i+2)\log L_{ii}.
\end{equation}
The collapsed log target for an IW covariance, up to constants in the proposal,
is therefore
\begin{equation}\label{eq:iw-log-cholesky-target}
 \ell(\bSigma,\ldots)-\frac{\nu+k+1}{2}\log|\bSigma|
 -\frac12\tr(\bm{\Psi}\bSigma^{-1})
 +\sum_{i=1}^k(k-i+2)\log L_{ii}.
\end{equation}
The constant $k\log2$ cancels for fixed dimension. For $k=1$ the full Jacobian
is $2L_{11}^2$, and for $k=2$ it is $4L_{11}^3L_{22}^2$.
These are covariance proposal coordinates; $\bL$ is not inserted into the fixed
observation mapping.

The Jacobian depends on the coordinate transformation and the measure used for
the prior density. Changing the IW degrees-of-freedom or scale convention changes
the density, not this transformation. Computing $\bSigma^{-1}$ for precision
assembly adds no further Jacobian. A covariance shared by several groups gets
one adjustment per estimated covariance matrix, not per group. Direct conjugate
IW Gibbs draws do not require a Metropolis adjustment.

\subsection{Alternative scale and correlation coordinates}
The following example uses different proposal coordinates and its own complete
Jacobian. It is not an additional correction to
\eqref{eq:iw-log-cholesky-target}.
For the two-coefficient case, write $v_j=\sigma_j^2=e^{\xi_j}$,
$\rho=\tanh\zeta$. The Jacobian from $(v_0,v_1,\rho)$ to the distinct covariance
entries $(\Sigma_{00},\Sigma_{11},\Sigma_{01})$ is $\sqrt{v_0v_1}$.
Consequently, an IW density specified on those covariance entries contributes
\begin{equation}\label{eq:iw-jacobian}
 \log p_{\Sigma}\{\bSigma(\xi_0,\xi_1,\zeta)\}
 +\tfrac32(\xi_0+\xi_1)+\log(1-\rho^2)
\end{equation}
to the transformed log-prior target. In contrast, separately specified densities
on $(v_0,v_1,\rho)$ contribute
\begin{equation}\label{eq:separation-jacobian}
 \log p_0(v_0)+\log p_1(v_1)+\log p_\rho(\rho)
 +\xi_0+\xi_1+\log(1-\rho^2).
\end{equation}
For a prior on a standard deviation, its variance density is
$p_v(v)=p_\sigma(\sqrt v)/(2\sqrt v)$ before applying the log-variance transform.
For $k>2$, independently transformed pairwise correlations do not guarantee a
positive-definite correlation matrix. A valid matrix parameterization and its
complete Jacobian are required.

\section{Augmented likelihoods and update order}
Conditional on augmentation variables, a Gaussian working likelihood can be
written using $\by^*$ and $\bW$, so that
\begin{equation}
 p(\by,\text{augmentation}\mid\bbeta,\bv,\ldots)
 \propto\exp\left[-\tfrac12
 (\by^*-\bo-\bX\bbeta-\bH\bv)\trans\bW
 (\by^*-\bo-\bX\bbeta-\bH\bv)\right].
\end{equation}
Proportionality here omits only terms constant in the parameters being updated;
it is not permission to omit augmentation-dependent normalizing terms when those
terms depend on a proposed parameter.

For binomial logit observations with trial counts $m_i$, let
$\eta_i=o_i+\bx_i\trans\bbeta+(\bH\bv)_i$, where
$\bx_i\trans$ denotes row $i$ of $\bX$. Then
\begin{equation}
 \omega_i\mid\cdots\sim\operatorname{PG}(m_i,\eta_i),\qquad
 y_i^*=\frac{y_i-m_i/2}{\omega_i},\qquad W_{ii}=\omega_i.
\end{equation}
For negative-binomial observations with mean $\exp(\eta_i)$ and size $r$,
\begin{equation}\label{eq:nb-working}
 \omega_i\mid\cdots\sim\operatorname{PG}(y_i+r,\eta_i-\log r),\qquad
 y_i^*=\frac{y_i-r}{2\omega_i}+\log r,\qquad W_{ii}=\omega_i.
\end{equation}
For probit observations, the working response is a unit-variance normal latent
response truncated above or below zero according to the observed outcome, and
$\bW=\bI$.

A jointly IID-collapsed iteration for an augmented model has the following order:
\begin{enumerate}
\item If NB size is estimated, update it using the ordinary NB likelihood with
the current full predictor and its prior.
\item Refresh augmentation variables conditional on the current full predictor
and the resulting size. Rebuild any working response/precision affected by size.
\item Update $\bbeta$ with all components of $\bv$ integrated out.
\item Update covariance/loading parameters using their marginal targets with
the augmentation held fixed.
\item Draw a compatible joint $\bv$ using the resulting parameters before the
next augmentation or size update that conditions on it.
\end{enumerate}
The Gaussian case omits augmentation and size updates. An ordinary NB size target
cannot be replaced by the covariance-only Gaussian working target while retaining
stale augmentation. Initialization and warmup must maintain the same conditioning
relationships. The separate-IID schedule instead uses explicit IID coefficients
in the coefficient and structured-covariance conditionals, with conditional IID
covariance updates as in \eqref{eq:conditional-covariance-target}.

\section{Vector-valued terms and multivariate observations}
\subsection{Four observation settings}
The same latent system covers the following observation layouts:
\begin{enumerate}
\item A univariate response with repeated observations sharing latent locations
or groups, and missing responses omitted from the likelihood.
\item A univariate response combining correlated coefficients within an IID or
structured term, such as an intercept and slope.
\item Multiple responses with correlated outcome-specific effects within a term,
including different observation locations and partially observed outcome vectors.
\item Multiple responses with several coefficients per outcome and within-term
correlation both within and across outcomes.
\end{enumerate}
Each actual observed scalar response supplies a row of $\by$, $\bX$, and $\bH$.
Observation rows, latent support points, outcomes, and coefficient components
are distinct indices. Repeated observations may reuse latent columns with
different covariate weights. Missing responses omit observation rows without
requiring removal of the corresponding latent support. No common set of
covariates or equal coefficient count across outcomes is required.

For example, two outcomes with a spatial intercept and slope have
\begin{equation}\label{eq:multivariate-svc-map}
 \bw(s)=\begin{pmatrix}w_{10}(s)\\w_{11}(s)\\w_{20}(s)\\w_{21}(s)\end{pmatrix},
 \qquad
 \bB_i=\begin{pmatrix}1&x_{i1}&0&0\\0&0&1&x_{i2}\end{pmatrix}.
\end{equation}
Their latent contribution is $\bB_i\bw(s_i)$. If only outcome 1 is observed,
only row 1 contributes to the likelihood. A $4\times4$ same-location covariance
can correlate intercepts and slopes within and across outcomes; the spatial
cross-covariance function additionally specifies dependence between locations.
For an IID term, analogous coefficient vectors are independent across groups.

If outcome $j$ has $k_j$ coefficients in a term, its total component dimension
is $d=\sum_j k_j$. An unrestricted covariance has $d(d+1)/2$ free parameters.
Six outcomes with intercepts and slopes therefore have $d=12$ and 78 covariance
parameters. The algebra imposes no small fixed dimension limit. Unrestricted
versus structured covariance, prior specifications, and proposal blocking are
modeling and computational choices at the resulting dimension.

\subsection{Correlated IID and spatial intercepts and slopes}
For a univariate response containing both terms, write
\begin{equation}
 y_i=\mu_i+\alpha_{0,g(i)}+x_i\alpha_{1,g(i)}
              +w_0(s_i)+x_iw_1(s_i)+\epsilon_i,
\end{equation}
where $\mu_i$ contains the fixed effects and offset. With
$\balpha_g\sim N_2(\bzero,\bSigma_{\mathrm{iid}})$ independently across groups,
the IID precision is $\bI_G\otimes\bSigma_{\mathrm{iid}}^{-1}$.
For an illustrative separable spatial specification,
\begin{equation}\label{eq:separable-svc}
 \Cov\{\bw(s),\bw(t)\}=r(s,t;\phi)\bSigma_{\mathrm{sp}},\qquad
 \bQ_w=\bR_s^{-1}\otimes\bSigma_{\mathrm{sp}}^{-1},
\end{equation}
where $\bw(s)=(w_0(s),w_1(s))\trans$, $r(s,s;\phi)=1$, and $\bR_s$ is the
spatial correlation matrix over distinct latent locations. The precision formula
uses location-major stacking of coefficient pairs. $\bR_s$ is distinct from the
observation-error covariance $\bR$.
The fixed design row places $[1\ \ x_i]$ in the group columns of $\bZ$ and in
the location columns of $\bA$. Thus the combined prior is
\begin{equation}\label{eq:iid-spatial-prior}
 \bQ=\blockdiag\{\bI_G\otimes\bSigma_{\mathrm{iid}}^{-1},
                 \bQ_w(\bSigma_{\mathrm{sp}},\phi)\}.
\end{equation}
These are the two prior blocks in \eqref{eq:M-blocks}. Each $2\times2$ covariance
has three free parameters shared across its groups or locations. The observation
cross-block $\bZ\trans\bW\bA$ creates posterior dependence between the terms;
it introduces no additional prior cross-covariance parameter between them.

The shared spatial correlation in \eqref{eq:separable-svc} is an example, not
a general requirement. Different component kernels require a valid joint
cross-covariance/precision construction, such as the one below. A same-location
covariance matrix alone does not specify the complete spatial model. Sparse
spatial formulations supply their appropriate precision construction; these
identities do not require explicitly inverting a large spatial covariance.

\subsection{Alternative covariance factorizations}
A linear model of coregionalization writes a response-specific spatial vector as
\begin{equation}\label{eq:lmc}
 \bw(s)=\bC\bu(s),\qquad
 \Cov\{\bu(s),\bu(t)\}=\diag\{r_1(s,t),\ldots,r_K(s,t)\}.
\end{equation}
The component processes are independent; each can have its own spatial kernel.
The resulting cross-covariance is
\begin{equation}
 \Cov\{\bw(s),\bw(t)\}=\bC\diag\{r_k(s,t)\}\bC\trans.
\end{equation}
Unknown entries of $\bC$ enter the observation mapping when the latent variables
are $\bu$. Alternatively, use $\bw$ as the latent vector and carry its joint
dependence in the prior. A reduced-rank loading construction may give $\bw$ a
singular prior; retaining $\bu$ then avoids applying the positive-definite
precision formulas to that singular covariance. These constructions connect
directly to multivariate process modeling and marginalization in \cite{Finley2008}.

Spatially varying loadings $\bC(s)$ give
$\bC(s)\diag\{r_k(s,t)\}\bC(t)\trans$ \cite{Gelfand2004}.
If the loadings themselves have a stochastic-process prior, the model is Gaussian
in $\bu$ conditional on those loadings. It is generally not jointly Gaussian in
both sets of unknowns because they enter multiplicatively. The Gaussian collapse
then conditions on the loading realization; it does not automatically integrate
both loading and component processes analytically.

Conditional multivariate models can instead give dependence directly to the
prior precision. For example, writing $\bw=\bB\bw+\bm{e}$ with strictly block
lower-triangular $\bB$ and innovation precision $\bLambda$ gives
\begin{equation}
 \bQ_w=(\bI-\bB)\trans\bLambda(\bI-\bB).
\end{equation}
This includes the conditional construction used by multivariate DAGAR models
\cite{Gao2022}. Such models express cross-component dependence directly in the
prior precision while keeping the observation mapping fixed.

\subsection{Known misalignment is an observation-selection problem}
Suppose response $j$ is observed only on locations $S_j$. Stack the actual
observations and let $\bS_{\mathrm{obs}}$ select the relevant response--location
pairs from a latent construction. For the selected centered representation, let
$\bH_{\mathrm{full}}$ contain the fixed selection and covariate weights for
the full observation layout. Then
\begin{equation}\label{eq:misalignment}
 \bH_{\mathrm{obs}}=\bS_{\mathrm{obs}}\bH_{\mathrm{full}}.
\end{equation}
For known locations and observation patterns, $\bS_{\mathrm{obs}}$ is fixed.
The row selections can also be constructed directly without storing unobserved
rows. Misalignment alone does not make the mapping stochastic. In an alternative
loading representation, $\bH_{\mathrm{full}}$ may depend on loading parameters.
A model with uncertain locations or unknown alignment would introduce additional
unknowns and is a different specification.
Response-specific designs and unequal observation counts fit this framework
\cite{Finley2014}; no fabricated observed outcomes are needed to create a complete
rectangular data array. Missing covariates require separate treatment.

Independent realizations within stands or ecoregions can be represented by prior
blocks for each parent. Shared versus parent-specific covariance parameters are
a separate modeling choice. Cross-response process covariance is also distinct
from cross-response observation-error covariance. The latter requires appropriate
blocks in $\bR$, rather than merely changing the latent loadings.

\section{Sparse computation and verification requirements}
The joint sparse implementation stores the coefficients in the order
$(\bw\trans,\balpha\trans)\trans$, retaining the existing offsets for structured
terms. This is a fixed permutation of the order used in \eqref{eq:M-blocks};
the observation mapping and prior precision are permuted together. It changes
neither the marginal likelihood nor the conditional distribution. The
structured-effect output dimension remains distinct from the joint system
dimension. Conditional structured-effect recovery given stored IID draws uses
the separate system, whereas joint recovery draws both coefficient blocks.

The same-group covariance in \eqref{eq:iid-crosscov} does not require additional
design entries. Nevertheless, including group nodes in the joint precision
\eqref{eq:M-blocks} can connect many structured nodes through $\bZ\trans\bW\bA$.
These connections can create Cholesky fill. A fill-reducing ordering can help,
but sparse input alone does not guarantee a sparse factor. Crossed IID terms can
also create fill without any structured process. Conversely, the separate
strategy can incur a dense IID Schur complement \eqref{eq:schur}.

The two representations and sampler strategies should be assessed using the
number of nonzeros in $\bM$ and its factor, memory, factorization and recovery
cost, and effective samples per second. Factorization patterns should be reused
when structural sparsity is fixed; changing numerical weights or sampled
loadings still requires the relevant numerical updates. Standardizing a CAR or
DAGAR scale does not generally make every node's marginal variance one.

Small independent numerical references should verify:
\begin{enumerate}
\item the intercept--slope covariance and the centered/standardized equivalence;
\item determinant and quadratic-form identities against a directly formed small
observation covariance;
\item the IID Schur complement and joint conditional moments;
\item transformed covariance-prior Jacobians and scalar special cases;
\item unchanged observation covariance when known misalignment is expressed by
selection rather than by direct row construction;
\item vector-valued terms covering all four observation layouts, including
different covariates by outcome, repeated measurements, and partially observed
outcome vectors;
\item restoration of all affected matrices and caches after rejected proposals.
\end{enumerate}
Explicit matrix inverses in such small reference calculations are mathematical
checks; the fitting implementation should use factors, solves, and small-block
operations.

\begin{thebibliography}{9}
\bibitem{Banerjee2008}
Banerjee, S., Gelfand, A. E., Finley, A. O., and Sang, H. (2008).
Gaussian predictive process models for large spatial data sets.
\emph{Journal of the Royal Statistical Society: Series B}, 70, 825--848.
\url{https://doi.org/10.1111/j.1467-9868.2008.00663.x}.

\bibitem{Finley2008}
Finley, A. O., Banerjee, S., Ek, A. R., and McRoberts, R. E. (2008).
Bayesian multivariate process modeling for prediction of forest attributes.
\emph{Journal of Agricultural, Biological, and Environmental Statistics},
13, 60--83. \url{https://doi.org/10.1198/108571108X273160}.

\bibitem{Gelfand2004}
Gelfand, A. E., Schmidt, A. M., Banerjee, S., and Sirmans, C. F. (2004).
Nonstationary multivariate process modeling through spatially varying
coregionalization. \emph{TEST}, 13, 263--312.

\bibitem{Finley2014}
Finley, A. O., Banerjee, S., and Cook, B. D. (2014).
Bayesian hierarchical models for spatially misaligned data in R.
\emph{Methods in Ecology and Evolution}, 5, 514--523.
\url{https://doi.org/10.1111/2041-210X.12189}.

\bibitem{Gao2022}
Gao, L., Datta, A., and Banerjee, S. (2022).
Hierarchical multivariate directed acyclic graph autoregressive models for spatial
diseases mapping. \emph{Statistics in Medicine}, 41, 3057--3075.
\url{https://doi.org/10.1002/sim.9404}.
\end{thebibliography}
\end{document}
